\clearpage

\renewcommand{\bibname}{\xiaoyi\bf\vspace{5pt} References}
\addcontentsline{toc}{chapter}{References}
\setstretch{1.25}

\begin{thebibliography}{99}
    \setlength{\itemindent}{-5pt}

    \bibitem{aapm} American Association of Physicists in Medicine. Low Dose CT Grand Challenge, 2016. Available: \url{https://www.aapm.org/GrandChallenge/LowDoseCT/}.

    \bibitem{redcnn} H. Chen, Y. Zhang, M. K. Kalra, F. Lin, Y. Chen, P. Liao, J. Zhou, and G. Wang, ``Low-dose CT with a residual encoder-decoder convolutional neural network,'' \textit{IEEE Transactions on Medical Imaging}, vol. 36, no. 12, pp. 2524--2535, 2017.

    \bibitem{ctformer} D. Wang, F. Fan, Z. Wu, R. Liu, F. Wang, and H. Yu, ``CTformer: Convolution-free token2token dilated vision Transformer for low-dose CT denoising,'' \textit{Physics in Medicine and Biology}, vol. 68, no. 6, 2023.

    \bibitem{restormer} S. W. Zamir, A. Arora, S. Khan, M. Hayat, F. S. Khan, M.-H. Yang, and L. Shao, ``Restormer: Efficient Transformer for high-resolution image restoration,'' in \textit{Proceedings of the IEEE/CVF Conference on Computer Vision and Pattern Recognition}, 2022, pp. 5728--5739.

    \bibitem{swinir} J. Liang, J. Cao, G. Sun, K. Zhang, L. Van Gool, and R. Timofte, ``SwinIR: Image restoration using Swin Transformer,'' in \textit{Proceedings of the IEEE/CVF International Conference on Computer Vision Workshops}, 2021, pp. 1833--1844.

    \bibitem{unet} O. Ronneberger, P. Fischer, and T. Brox, ``U-Net: Convolutional networks for biomedical image segmentation,'' in \textit{Medical Image Computing and Computer-Assisted Intervention}, 2015, pp. 234--241.

    \bibitem{attention} A. Vaswani et al., ``Attention is all you need,'' in \textit{Advances in Neural Information Processing Systems}, 2017, pp. 5998--6008.

    \bibitem{vit} A. Dosovitskiy et al., ``An image is worth 16x16 words: Transformers for image recognition at scale,'' in \textit{International Conference on Learning Representations}, 2021.

    \bibitem{ssim} Z. Wang, A. C. Bovik, H. R. Sheikh, and E. P. Simoncelli, ``Image quality assessment: From error visibility to structural similarity,'' \textit{IEEE Transactions on Image Processing}, vol. 13, no. 4, pp. 600--612, 2004.

    \bibitem{noise2noise} J. Lehtinen et al., ``Noise2Noise: Learning image restoration without clean data,'' in \textit{Proceedings of the International Conference on Machine Learning}, 2018.

    \bibitem{noise2self} J. Batson and L. Royer, ``Noise2Self: Blind denoising by self-supervision,'' in \textit{Proceedings of the International Conference on Machine Learning}, 2019.

    \bibitem{noise2void} A. Krull, T.-O. Buchholz, and F. Jug, ``Noise2Void: Learning denoising from single noisy images,'' in \textit{Proceedings of the IEEE/CVF Conference on Computer Vision and Pattern Recognition}, 2019.

    \bibitem{framelet} E. Kang, W. Chang, J. Yoo, and J. C. Ye, ``Deep convolutional framelet denoising for low-dose CT via wavelet residual network,'' \textit{IEEE Transactions on Medical Imaging}, vol. 37, no. 6, pp. 1358--1369, 2018.

    \bibitem{fdk} L. A. Feldkamp, L. C. Davis, and J. W. Kress, ``Practical cone-beam algorithm,'' \textit{Journal of the Optical Society of America A}, vol. 1, no. 6, pp. 612--619, 1984.

    \bibitem{mbir} J.-B. Thibault, K. D. Sauer, C. A. Bouman, and J. Hsieh, ``A three-dimensional statistical approach to improved image quality for multislice helical CT,'' \textit{Medical Physics}, vol. 34, no. 11, pp. 4526--4544, 2007.

\end{thebibliography}
\doublespacing
